% ============================================================
%  Chapter 1 -- Euclidean Spaces, 1.36-1.38
% ============================================================
\rsection{Euclidean Spaces}

\begin{signpost}[The last piece of equipment]
Three items, and the section exists for one sentence: Remark 1.38, which
observes that $\R^k$ is a metric space. Everything in Chapter 2 --- open sets,
compactness, connectedness --- is developed for metric spaces, and $\R^k$ is the
example that motivates all of it. This section supplies the distance function
and verifies the three properties Chapter 2 will demand.

\smallskip
Note what is \emph{not} claimed. Rudin gives $\R^k$ a vector space structure and
an inner product, but no multiplication of vectors by vectors --- $\R^k$ is not
a field for $k > 1$. The only exception is $k = 2$, which coincides with $\C$;
1.38 remarks on the coincidence, and it is worth keeping straight that $\R^2$
and $\C$ are the same metric space but only $\C$ has a product.
\end{signpost}

\ritem{1.36}{Definitions}
For each positive integer $k$, let $\R^k$ be the set of all ordered $k$-tuples
\[ \mathbf{x} = (x_1, x_2, \dots, x_k), \]
where $x_1, \dots, x_k$ are real numbers, called the \emph{coordinates} of
$\mathbf{x}$. The elements of $\R^k$ are called \emph{points}, or
\emph{vectors}, especially when $k > 1$. We shall denote vectors by boldfaced
letters. If $\mathbf{y} = (y_1,\dots,y_k)$ and if $\alpha$ is a real number, put
\begin{align*}
  \mathbf{x}+\mathbf{y} &= (x_1+y_1, \dots, x_k+y_k)\\
  \alpha\mathbf{x} &= (\alpha x_1, \dots, \alpha x_k)
\end{align*}
so that $\mathbf{x}+\mathbf{y} \in \R^k$ and $\alpha\mathbf{x} \in \R^k$. This
defines addition of vectors, as well as multiplication of a vector by a real
number (a \emph{scalar}). These two operations satisfy the commutative,
associative, and distributive laws (the proof is trivial, in view of the
analogous laws for the real numbers) and make $\R^k$ into a \emph{vector space
over the real field}. The zero element of $\R^k$ (sometimes called the
\emph{origin} or the \emph{null vector}) is the point $\mathbf{0}$, all of whose
coordinates are $0$.

We also define the so-called ``inner product'' (or scalar product) of
$\mathbf{x}$ and $\mathbf{y}$ by
\[ \mathbf{x}\cdot\mathbf{y} = \sum_{i=1}^k x_iy_i \]
and the \emph{norm} of $\mathbf{x}$ by
\[ |\mathbf{x}| = (\mathbf{x}\cdot\mathbf{x})^{1/2}
   = \left(\sum_{i=1}^k x_i^2\right)^{1/2}. \]

The structure now defined (the vector space $\R^k$ with the above inner product
and norm) is called \emph{euclidean $k$-space}.\kn{Three separate structures are
being stacked here, and later chapters use them independently: a \emph{vector
space} (addition and scaling, used throughout Chapter 9), an \emph{inner
product} (used for orthogonality and in Chapter 11), and a \emph{norm} (used to
make a metric, which is all Chapter 2 needs). The norm is derived from the inner
product; the inner product is not derived from anything.}

\ritem{1.37}{Theorem}
\emph{Suppose $\mathbf{x},\mathbf{y},\mathbf{z} \in \R^k$, and $\alpha$ is real.
Then}
\begin{enumerate}[label=(\alph*), leftmargin=2.2em, itemsep=1pt, topsep=3pt]
  \item \emph{$|\mathbf{x}| \ge 0$;}
  \item \emph{$|\mathbf{x}| = 0$ if and only if $\mathbf{x} = \mathbf{0}$;}
  \item \emph{$|\alpha\mathbf{x}| = |\alpha|\,|\mathbf{x}|$;}
  \item \emph{$|\mathbf{x}\cdot\mathbf{y}| \le |\mathbf{x}|\,|\mathbf{y}|$;}
  \item \emph{$|\mathbf{x}+\mathbf{y}| \le |\mathbf{x}|+|\mathbf{y}|$;}
  \item \emph{$|\mathbf{x}-\mathbf{z}| \le |\mathbf{x}-\mathbf{y}|
        + |\mathbf{y}-\mathbf{z}|$.}
\end{enumerate}

\begin{proof}
(a), (b), and (c) are obvious, and (d) is an immediate consequence of the
Schwarz inequality.\kn{Take 1.35 with all $a_j, b_j$ real; conjugation does
nothing, and the statement becomes
$(\sum x_iy_i)^2 \le \sum x_i^2 \sum y_i^2$, which is (d) squared. This is the
only place the previous section's main theorem is used, and it is the reason
that section had to come first.} By (d) we have
\begin{align*}
  |\mathbf{x}+\mathbf{y}|^2
    &= (\mathbf{x}+\mathbf{y})\cdot(\mathbf{x}+\mathbf{y})\\
    &= \mathbf{x}\cdot\mathbf{x} + 2\,\mathbf{x}\cdot\mathbf{y}
       + \mathbf{y}\cdot\mathbf{y}\\
    &\le |\mathbf{x}|^2 + 2|\mathbf{x}|\,|\mathbf{y}| + |\mathbf{y}|^2\\
    &= (|\mathbf{x}|+|\mathbf{y}|)^2,
\end{align*}
so that (e) is proved. Finally (f) follows from (e) if we replace $\mathbf{x}$
by $\mathbf{x}-\mathbf{y}$ and $\mathbf{y}$ by $\mathbf{y}-\mathbf{z}$.
\end{proof}

\begin{deeper}[The geometry of 1.37(d) and (e)]
Parts (d) and (e) are the two facts of Euclidean geometry that survive into
arbitrary dimension, and each has a picture that is also its proof.

\medskip
\textbf{(d) Schwarz is a statement about projection.} Assume
$\mathbf{y} \ne \mathbf{0}$ --- otherwise both sides of (d) vanish and there
is nothing to prove. Drop a perpendicular from $\mathbf{x}$ onto the line
through $\mathbf{y}$. The foot sits at $t\mathbf{y}$ with
$t = (\mathbf{x}\cdot\mathbf{y})/|\mathbf{y}|^2$, and the leftover piece
$\mathbf{x}-t\mathbf{y}$ is perpendicular to $\mathbf{y}$.

\medskip
\begin{center}
\begin{tikzpicture}[x=1mm, y=1mm]
  \draw[thin] (-4,0) -- (56,0);
  % the projection, drawn heavy along the line
  \draw[seg] (0,0) -- (26,0);
  \draw[vec] (0,0) -- (44,0);  \node[mlbl, anchor=south] at (44,1.4) {$\mathbf{y}$};
  \draw[vec] (0,0) -- (26,17);  \node[mlbl, anchor=south east] at (26,17) {$\mathbf{x}$};
  \draw[guide] (26,17) -- (26,0);
  \draw[thin] (22.5,0) -- (22.5,3.5) -- (26,3.5);
  \node[ptfill, minimum size=2.6pt] at (26,0) {};
  \node[mlbl, anchor=west] at (27.5,2.5) {$t\mathbf{y}$};
  \node[lbl, anchor=west] at (28,11) {$\mathbf{x}-t\mathbf{y}$};
  % the projection length, safely below the line
  \draw[thin, {Stealth[length=3pt]}-{Stealth[length=3pt]}] (0,-4) -- (26,-4);
  \node[lbl, anchor=north] at (13,-5.5)
    {this shadow has length $|\mathbf{x}\cdot\mathbf{y}|/|\mathbf{y}|$};
\end{tikzpicture}
\end{center}
\figcap{The shadow of $\mathbf{x}$ on $\mathbf{y}$ cannot be longer than
$\mathbf{x}$ itself.}

\smallskip
Now compute the length of that leftover piece:
\[
  0 \;\le\; |\mathbf{x}-t\mathbf{y}|^2
    \;=\; |\mathbf{x}|^2 - \frac{(\mathbf{x}\cdot\mathbf{y})^2}{|\mathbf{y}|^2},
\]
and rearranging gives exactly
$(\mathbf{x}\cdot\mathbf{y})^2 \le |\mathbf{x}|^2|\mathbf{y}|^2$. That is
Schwarz, and no angle was needed --- only that a squared length is not
negative. Equality holds precisely when the leftover piece vanishes, that is
when $\mathbf{x}$ lies on the line through $\mathbf{y}$: the linear dependence
of Exercise 1.15.

\smallskip
This is also Rudin's proof of 1.35 in disguise. His $\sum|Ba_j - Cb_j|^2 \ge 0$
is the same ``a squared length is not negative,'' with $Ba_j - Cb_j$ playing
the part of the leftover piece --- and his separate treatment of $B=0$
is exactly the $\mathbf{y}=\mathbf{0}$ case set aside above.

\smallskip
A third way to say it, if you prefer algebra to pictures: the function
$q(t) = |\mathbf{x}-t\mathbf{y}|^2$ is a quadratic in $t$ that is never
negative, so its discriminant cannot be positive, and
$4(\mathbf{x}\cdot\mathbf{y})^2 - 4|\mathbf{y}|^2|\mathbf{x}|^2 \le 0$. The
three proofs are one proof: the projection is where $q$ attains its minimum,
and the discriminant condition says that minimum is not below zero.

\medskip
\textbf{(e) The triangle inequality is a triangle.} Lay $\mathbf{y}$ head to
tail after $\mathbf{x}$; the direct route to $\mathbf{x}+\mathbf{y}$ is no
longer than going via the corner.

\medskip
\begin{center}
\begin{tikzpicture}[x=1mm, y=1mm]
  \draw[vec] (0,0) -- (30,4);   \node[mlbl, anchor=north] at (15,2.6) {$\mathbf{x}$};
  \draw[vec] (30,4) -- (44,20); \node[mlbl, anchor=west] at (38,11) {$\mathbf{y}$};
  \draw[vec] (0,0) -- (44,20);  \node[mlbl, anchor=south east] at (24,13.5) {$\mathbf{x}+\mathbf{y}$};
  \node[ptfill, minimum size=2.6pt] at (0,0) {};
  \node[ptfill, minimum size=2.6pt] at (30,4) {};
  \node[ptfill, minimum size=2.6pt] at (44,20) {};
  \node[lbl, anchor=west] at (50,10) {direct route $\le$};
  \node[lbl, anchor=west] at (50,6.5) {route via the corner};
\end{tikzpicture}
\end{center}
\figcap{Part (f) is the same picture read as three \emph{points}: going from
$\mathbf{x}$ to $\mathbf{z}$ directly beats detouring through $\mathbf{y}$.}

\smallskip
\textbf{But here the picture is not a proof, and it is worth being clear why.}
For Schwarz the figure could be turned into an argument, because every step
was an identity about inner products. Not so here. ``The straight path is
shortest'' is a theorem \emph{of} Euclidean geometry, and in Rudin's
development Euclidean geometry does not yet exist --- 1.36 has just
\emph{defined} distance by a formula, and 1.37 is where the familiar
properties of that formula are established for the first time. Appealing to
the picture would assume what is being proved. The figure shows you what
1.37(e) means; the algebra, via Schwarz, is what makes it true.
\end{deeper}

\begin{unpack}[Part (f) is the triangle inequality you will actually use]
Parts (e) and (f) say the same thing twice, and it is worth seeing why both are
stated. Part (e) is about \emph{vectors}: the length of a sum is at most the sum
of the lengths. Part (f) is about \emph{points}: the distance from $\mathbf{x}$
to $\mathbf{z}$ is at most the distance via any detour through $\mathbf{y}$.

The substitution converting one into the other is worth doing once by hand.
Setting $\mathbf{u} = \mathbf{x}-\mathbf{y}$ and
$\mathbf{v} = \mathbf{y}-\mathbf{z}$, we have
$\mathbf{u}+\mathbf{v} = \mathbf{x}-\mathbf{z}$, so (e) reads
$|\mathbf{x}-\mathbf{z}| \le |\mathbf{x}-\mathbf{y}|+|\mathbf{y}-\mathbf{z}|$.

Form (f) is the one that generalises. When Chapter 2 defines a metric space
there is no addition available, so vectors cannot be summed and (e) has no
meaning --- but distances between three points still make sense, and (f) becomes
axiom (c) of Definition 2.15.
\end{unpack}

\ritem{1.38}{Remarks}
Theorem 1.37(a), (b), and (f) will allow us (see Chap.\ 2) to regard $\R^k$ as a
\emph{metric space}.\kd{This is the entire purpose of the section, stated in one
sentence and easy to read past. Compare the three parts he names with Definition
2.15: (a) and (b) give $d(x,y)>0$ for $x \ne y$ and $d(x,x)=0$; (f) is the
triangle inequality. Symmetry, the remaining axiom, follows from (c) with
$\alpha = -1$. Rudin has just verified, in advance, that $\R^k$ is an example of
the structure Chapter 2 will study abstractly.}

$\R^1$ (the set of all real numbers) is usually called \emph{the line}, or
\emph{the real line}. Likewise, $\R^2$ is called \emph{the plane}, or \emph{the
complex plane} (compare Definitions 1.24 and 1.36). In these two cases, the norm
is just the absolute value of the corresponding real or complex number.

\begin{pitfall}[$\R^2$ and $\C$ are the same space but not the same object]
As metric spaces they are identical: $(a,b) \in \R^2$ and $a+bi \in \C$ have the
same norm $\sqrt{a^2+b^2}$, so every statement about convergence, open sets or
compactness transfers word for word.

They differ in algebra. $\C$ has a multiplication making it a field (1.25);
$\R^2$ has only the inner product, which produces a \emph{scalar}, not a vector.
So $zw$ makes sense for $z,w \in \C$ and $\mathbf{x}\mathbf{y}$ does not for
$\mathbf{x},\mathbf{y} \in \R^2$.

The distinction matters when a theorem is stated for one and you want it for the
other. Results resting only on the metric --- most of Chapter 2, and the
convergence theorems of Chapter 3 --- carry across freely. Results using
multiplication or division, such as the whole theory of power series in Chapter
8, belong to $\C$ alone.
\end{pitfall}

\begin{recap}[The chapter proper, in one paragraph]
$\Q$ is unfit for analysis because a bounded set can fail to have a least upper
bound (1.1). Name that failure (1.8--1.10), develop just enough algebra to state
what $\R$ must be (1.12--1.18), assert a field free of the failure (1.19), and
the consequences follow at once: infima exist too (1.11), $\N$ is unbounded and
$\Q$ is dense (1.20), and every positive real has every root (1.21). The
remaining sections build the objects later chapters will work in --- the
extended reals for unbounded suprema, $\C$ for Chapters 3 and 8, $\R^k$ for
Chapter 2 onward --- and prove one inequality of lasting importance, Schwarz's
(1.35). Nothing in the rest of the book uses more about $\R$ than 1.19. When a
later proof says ``by completeness,'' that is what it means.
\end{recap}
